\chapter{Hyperacusis}

\section*{Definition}
Hyperacusis is an abnormal intolerance to everyday sounds that are not considered loud by the general population, often accompanied by pain, annoyance, or fear in response to moderate-intensity sounds (typically 50--65 dB).

\section*{What Happens in Your Body}
Dysregulation of central auditory gain mechanisms, particularly involving the olivocochlear efferent system and serotonergic pathways, leads to amplified neural response in the auditory cortex. Peripheral triggers include cochlear hair cell damage and stapedius muscle weakness.

\section*{Causes}
\begin{itemize}
  \item \textbf{Noise-induced hearing loss:} Chronic noise exposure causing cochlear damage
  \item \textbf{Facial nerve palsy (Bell palsy):} Stapedius muscle paralysis reduces acoustic reflex dampening
  \item \textbf{Meniere disease:} Endolymphatic hydrops with fluctuating hearing sensitivity
  \item \textbf{Superior semicircular canal dehiscence:} Third window effect amplifying bone-conducted sound
  \item \textbf{Migraine:} Central sensitization with phonophobia (often overlaps with hyperacusis)
  \item \textbf{Autism spectrum disorder:} Sensory processing differences with sound hypersensitivity
  \item \textbf{Post-traumatic stress disorder:} Hypervigilance and heightened acoustic startle response
\end{itemize}

\section*{When to See a Doctor}
See your doctor if:
\begin{itemize}
  \item Sound intolerance interferes with work, school, or social function
  \item Associated with hearing loss, tinnitus, or vertigo
  \item Acute onset after head trauma or barotrauma (suspect perilymph fistula or stapes injury)
  \item Unilateral hyperacusis with facial weakness (suspect acoustic neuroma or parotid tumor)
  \item Severe phonophobia or misophonia causing avoidance behaviors
\end{itemize}

\section*{Related Symptoms}
\begin{itemize}
  \item Tinnitus (coexisting in majority of cases)
  \item Hearing loss (may be present with associated sensorineural loss)
  \item Ear pain or aural fullness
  \item Vertigo or dizziness
  \item Phonophobia (fear of sound) and misophonia (hatred of specific sounds)
\end{itemize}
\textbf{Important:} Hyperacusis distinct from phonophobia (a migraine symptom) in that the discomfort is loudness-based rather than specifically fear-based; both can co-occur.

\section*{Self-Care / Home Management}
\begin{itemize}
  \item Sound desensitization therapy: gradual, graded exposure to recorded environmental sounds at low volume
  \item Earplugs should be used sparingly (only in unavoidable loud environments), as overprotection worsens central sensitization
  \item White noise generators or wearable sound generators at a volume below the discomfort threshold
  \item Relaxation techniques (deep breathing, mindfulness) to reduce anxiety associated with sound exposure
  \item Avoidance of complete silence, which may paradoxically increase auditory gain
\end{itemize}

\section*{How Your Doctor Will Evaluate You}
\begin{itemize}
  \item Comprehensive audiometry including pure-tone thresholds and loudness discomfort levels (LDLs)
  \item LDL typically below 80--90 dB HL in hyperacusis (normal LDL > 100 dB HL)
  \item Acoustic reflex testing to assess stapedius function
  \item MRI with gadolinium if unilateral hyperacusis and asymmetric audiometry
  \item Otoacoustic emissions to evaluate cochlear outer hair cell function
  \item Tinnitus Handicap Inventory and Hyperacusis Questionnaire for severity assessment
\end{itemize}

\section*{Treatment Options}
\begin{itemize}
  \item Cognitive-behavioral therapy (CBT) to address the emotional and behavioral response to sound
  \item Tinnitus retraining therapy (TRT) using sound generators at the mixing point
  \item Desensitization program with environmental sound enrichment (not silence)
  \item Treatment of underlying condition: corticosteroids for Bell palsy, diuretic + low-sodium diet for Meniere, migraine prophylaxis
  \item Hearing aids with compression settings for hyperacusis with coexisting hearing loss
  \item Surgical repair for superior canal dehiscence or perilymph fistula if identified
\end{itemize}

\section*{References}

\begin{thebibliography}{99}
\bibitem{harrison-hearing2} Longo DL, et al. \emph{Harrison's Principles of Internal Medicine}. 21st ed. McGraw-Hill; 2022. Chapter 44: Hearing Loss and Hyperacusis.
\bibitem{uptodate-hyperac} Baguley DM, et al. Hyperacusis in adults. In: UpToDate, Post TW (Ed), UpToDate, Waltham, MA; 2025.
\bibitem{tyler} Tyler RS, et al. Hyperacusis: a review. \emph{Am J Audiol}. 2023;32(1):1-17.
\bibitem{baguley} Baguley DM, et al. Hyperacusis and its clinical management. \emph{Trends Hear}. 2023;27:23312165231168450.
\bibitem{bartnik} Bartnik G, et al. Management of hyperacusis: a systematic review. \emph{J Am Acad Audiol}. 2023;34(3-4):79-92.
\bibitem{aroz} Aroz F, et al. Sound tolerance disorders: a clinical approach. \emph{Ear Hear}. 2023;44(4):719-731.
\end{thebibliography}

\clearpage
